
\documentclass{article}
\usepackage{colortbl}
\usepackage{makecell}
\usepackage{multirow}
\usepackage{supertabular}

\begin{document}

\newcounter{utterance}

\centering \large Interaction Transcript for game `eqbench', experiment `eq\_bench', episode 45 with reflexion{-}Qwen3.5{-}27B.
\vspace{24pt}

{ \footnotesize  \setcounter{utterance}{1}
\setlength{\tabcolsep}{0pt}
\begin{supertabular}{c@{$\;$}|p{.15\linewidth}@{}p{.15\linewidth}p{.15\linewidth}p{.15\linewidth}p{.15\linewidth}p{.15\linewidth}}
    \# & \multicolumn{2}{c}{Player} && \multicolumn{2}{c}{Game Master} \\
    \hline

    \theutterance \stepcounter{utterance}  
    & & & \multicolumn{4}{p{0.6\linewidth}}{
        \cellcolor[rgb]{0.9,0.9,0.9}{
            \makecell[{{p{\linewidth}}}]{
                \texttt{\tiny{[P1$\langle$GM]}}
                \texttt{You have attempted this same task before and did not succeed. These are the plans you wrote after those attempts.} \\
\texttt{Trial \#1: New plan: My previous attempt failed because I assigned a high score to "Defensive" (9) and "Resentful" (7) without sufficiently distinguishing the nuance of the final exchange. Brandon's final line ("And yet, you sell your installations, don't you?") is a direct challenge to Cecilia's moral high ground, which typically triggers a defensive reaction, but her immediate rebuttal ("I sell my work to survive, not to profit") is a rationalization rather than an emotional outburst. I likely overestimated the "Defensive" score because I focused on the conflict rather than her specific response style. Additionally, I underestimated "Guilt" (2); while she denies it, the fact that she feels the need to justify her survival vs. profit distinction suggests a lingering, low{-}level guilt or cognitive dissonance that I missed. Next time, I will analyze the *specific* rebuttal to the final accusation to calibrate the "Defensive" score more accurately (likely lowering it if the response is calm) and re{-}evaluate "Guilt" based on the necessity of the justification. I will also ensure the scores reflect a hierarchy where the primary emotion (likely Resentful or Defensive) is distinct from secondary ones, avoiding a flat distribution of high scores. Specifically, I will look for cues of "moral superiority" to boost "Proud" if she doubles down, or "shame" to boost "Guilt" if she hesitates, rather than defaulting to high defensiveness.} \\
\\ 
\texttt{Trial \#2: New plan: My previous attempt failed because I assigned high scores to multiple conflicting emotions (Defensive: 7, Resentful: 6, Proud: 4) without identifying the single dominant emotional state triggered by the final exchange. Brandon's accusation ("And yet, you sell your installations...") directly attacks Cecilia's moral consistency, and her immediate, sharp rebuttal ("I sell my work to survive, not to profit") is a classic defensive maneuver to protect her self{-}image against a perceived hypocrisy charge. I incorrectly diluted the intensity by spreading scores across "Resentful" and "Proud," which are less immediate reactions to a direct personal attack than "Defensive." Next time, I will prioritize the emotion that directly addresses the specific nature of the final provocation. Since the provocation is an attack on her integrity, I will assign the highest score (8{-}9) to "Defensive" as the primary reaction. I will significantly lower "Resentful" (to 2{-}3) as she is countering rather than stewing, and "Proud" (to 1{-}2) as she is justifying rather than boasting. I will keep "Guilt" low (1{-}2) as she explicitly denies the premise of the accusation rather than accepting it. The key is to ensure one score is distinctly higher than the others to reflect the primary emotional driver of the specific dialogue turn.} \\
\\ 
\texttt{Trial \#3: New plan: My previous attempts failed because I misidentified the primary emotional driver of the final exchange. While Brandon's accusation ("And yet, you sell your installations...") is an attack, Cecilia's response ("I sell my work to survive, not to profit") is not a defensive outburst or a sign of guilt; it is a firm, principled correction that reinforces her moral superiority over his commodification of art. By assigning a high score to "Defensive" (9) and "Resentful" (4), I missed that her tone is one of conviction and disdain for his worldview, not fear or anger. The emotion that best fits her stance of correcting his "delusion" and asserting the "soul" of art against his "commodity" view is **Resentful** (at his reductionist view) or **Proud** (in her own integrity), but given the sharpness of the rebuttal to his hypocrisy, **Resentful** is the most accurate primary emotion, as she is rejecting his entire philosophy with a sense of moral indignation. "Defensive" should be low because she is not backing down; she is advancing her argument. "Guilt" should be 0 as she explicitly rejects the premise. "Proud" is secondary to the active rejection of his view. Next time, I will assign the highest score (8{-}9) to **Resentful** to capture her moral indignation at his hypocrisy, lower **Defensive** to 2{-}3 (as she is countering, not cowering), set **Guilt** to 0, and keep **Proud** moderate (4{-}5) as a supporting state. The key is to recognize that a calm, sharp rebuttal to a moral attack often signals resentment of the attacker's worldview rather than personal defensiveness.} \\
\\ 
\\ 
\\ 
\texttt{Your task is to predict the likely emotional responses of a character in this dialogue:} \\
\\ 
\texttt{Cecilia: You know, your words have power, Brandon. More than you might think.} \\
\texttt{Brandon: I'm well aware, Cecilia. It's a critic's job to wield them.} \\
\texttt{Cecilia: But do you understand the weight of them? The lives they can shatter?} \\
\texttt{Brandon: Art is not for the faint{-}hearted. If you can't handle the critique, you're in the wrong industry.} \\
\texttt{Cecilia: It's not about handling criticism, Brandon. It's about understanding the soul of the art. You dissect it like a cold, lifeless body on an autopsy table.} \\
\texttt{Brandon: And you, Cecilia, romanticize it to the point of delusion. Art is not some sacred, untouchable entity. It's a commodity, like everything else.} \\
\texttt{Cecilia: That's where you're wrong. Art is a mirror of the human soul, a reflection of our deepest fears and desires. It's not a product to be bought and sold.} \\
\texttt{Brandon: And yet, you sell your installations, don't you? How is that different?} \\
\texttt{Cecilia: I sell my work to survive, not to profit. There's a difference.} \\
\texttt{{[}End dialogue{]}} \\
\\ 
\texttt{At the end of this dialogue, Cecilia would feel...} \\
\texttt{Defensive} \\
\texttt{Guilt} \\
\texttt{Proud} \\
\texttt{Resentful} \\
\\ 
\texttt{Give each of these possible emotions a score from 0{-}10 for the relative intensity that they are likely to be feeling each.} \\
\\ 
\texttt{You must output in the following format, including headings (of course, you should give your own scores), with no additional commentary:} \\
\\ 
\texttt{Defensive: <score>} \\
\texttt{Guilt: <score>} \\
\texttt{Proud: <score>} \\
\texttt{Resentful: <score>} \\
\\ 
\\ 
\texttt{{[}End of answer{]}} \\
\\ 
\texttt{Remember: zero is a valid score, meaning they are likely not feeling that emotion. You must score at least one emotion > 0.} \\
\\ 
\texttt{Your answer:} \\
            }
        }
    }
     \\ \\

    \theutterance \stepcounter{utterance}  
    & \multicolumn{4}{p{0.6\linewidth}}{
        \cellcolor[rgb]{0.561,0.737,0.561}{
            \makecell[{{p{\linewidth}}}]{
                \texttt{\tiny{[P1$\rangle$GM]}}
                \texttt{Defensive: 3} \\
\texttt{Guilt: 0} \\
\texttt{Proud: 6} \\
\texttt{Resentful: 8} \\
            }
        }
    }
    & & \\ \\

    \theutterance \stepcounter{utterance}  
    & & & \multicolumn{2}{p{0.3\linewidth}}{
        \cellcolor[rgb]{0.9,0.9,0.9}{
            \makecell[{{p{\linewidth}}}]{
                \texttt{\tiny{[GM$|$GM]}}
                \texttt{Defensive: 3} \\
\texttt{Guilt: 0} \\
\texttt{Proud: 6} \\
\texttt{Resentful: 8} \\
            }
        }
    }
    & & \\ \\

    \theutterance \stepcounter{utterance}  
    & & & \multicolumn{2}{p{0.3\linewidth}}{
        \cellcolor[rgb]{0.9,0.9,0.9}{
            \makecell[{{p{\linewidth}}}]{
                \texttt{\tiny{[GM$|$GM]}}
                \texttt{Defensive: 8} \\
\texttt{Guilt: 0} \\
\texttt{Proud: 4} \\
\texttt{Resentful: 7} \\
            }
        }
    }
    & & \\ \\

    \theutterance \stepcounter{utterance}  
    & & & \multicolumn{2}{p{0.3\linewidth}}{
        \cellcolor[rgb]{0.9,0.9,0.9}{
            \makecell[{{p{\linewidth}}}]{
                \texttt{\tiny{[GM$|$GM]}}
                \texttt{game\_result = LOSE} \\
            }
        }
    }
    & & \\ \\

    \theutterance \stepcounter{utterance}  
    & & \multicolumn{4}{p{0.6\linewidth}}{
        \cellcolor[rgb]{0.95,0.95,0.95}{
            \makecell[{{p{\linewidth}}}]{
                \texttt{\tiny{[P1$\leftrightarrow$P]}}
                \texttt{Post{-}episode Reflexion (not sent to GM):} \\
\texttt{New plan: My previous attempt failed because I misinterpreted the nuance of Cecilia's final rebuttal. While Brandon's accusation ("And yet, you sell your installations...") is a direct challenge, Cecilia's response ("I sell my work to survive, not to profit") is a calm, principled distinction rather than an emotional outburst. I incorrectly assigned a high score to "Resentful" (8) and "Proud" (6), assuming her firmness indicated moral indignation or superiority. However, the act of immediately justifying one's actions against a charge of hypocrisy is the definition of being **Defensive**. She is actively protecting her moral standing against his attack, not merely expressing disdain (Resentful) or boasting (Proud). The "survive vs. profit" distinction is a shield, not a sword. Next time, I will identify the *function* of the character's response: if they are explaining or justifying an action to counter an accusation, the primary emotion is "Defensive." I will assign the highest score (8{-}9) to **Defensive**, significantly lower **Resentful** (2{-}3) as she is engaging in logic rather than stewing in anger, lower **Proud** (1{-}2) as she is justifying necessity rather than celebrating achievement, and keep **Guilt** at 0 since she explicitly rejects the premise of the accusation. The key is to recognize that a rationalization in response to a personal attack is a defensive mechanism.} \\
            }
        }
    }
    & \\ \\

\end{supertabular}
}

\end{document}
